\scanpage{114}

CORNU\qquad Ph\qquad 16/06/78
\mspage{$\overline{\text{\uline{VI}}}$}

C4.Math

\figure{\begin{tikzpicture}[x=1cm,y=1cm]\draw (0,0) .. controls (2,.02) and (5,.05) .. (7.3,.1);\end{tikzpicture}}

Si l'on \unsure{associ} à chaque trièdre, une Couleur, on obtient ainsi un coloriage des arêtes de l'icosaèdre, tel que deux arêtes distinctes ayant la meme Couleur, \deleted{\unsure{ne} \unsure{se}} n'ont jamais de Sommets en Commun.

En relation avec tout ça, il aurait fallu parler ; du groupe des automorphismes de l'icosaèdre, dont le neutre et l'antipodisme Constituent le Centre.

\deleted{\unsure{ordre}}

\begin{center}
\uline{\unsure{Cornu}}.
\end{center}

\margin{113}
